\documentclass[10pt]{article}
\usepackage[T1]{fontenc}
\usepackage[utf8]{inputenc}
\usepackage[french]{babel}
\usepackage{newpxtext,newpxmath}
\usepackage[a4paper,top=1.5cm,bottom=1.35cm,left=1.6cm,right=1.6cm]{geometry}
\usepackage{microtype}
\usepackage{xcolor}
\usepackage{titlesec}
\usepackage[most]{tcolorbox}
\usepackage{enumitem}
\usepackage{array}
\usepackage{booktabs}
\usepackage{fancyhdr}

\definecolor{navy}{HTML}{1F3A5F}
\definecolor{grey}{HTML}{54606E}
\definecolor{green}{HTML}{2C6E49}
\definecolor{amber}{HTML}{9C6A16}
\definecolor{red}{HTML}{A6322E}
\definecolor{plum}{HTML}{7A4EA8}
\definecolor{boxgray}{HTML}{F3F5F8}
\definecolor{boxred}{HTML}{FCF2F1}
\definecolor{boxgreen}{HTML}{F1F7F3}
\definecolor{boxplum}{HTML}{F7F3FB}
\definecolor{boxcode}{HTML}{F7F8FA}
\definecolor{boxedge}{HTML}{D9E0EA}

\setlength{\parindent}{0pt}
\setlength{\parskip}{2.6pt}
\linespread{1.0}

\pagestyle{fancy}\fancyhf{}
\lhead{\small\color{grey}Scella}
\rhead{\small\color{grey}Le groupe le plus rachetable -- 11 août 2026}
\cfoot{\small\color{grey}\thepage}
\renewcommand{\headrulewidth}{0.3pt}

\titleformat{\section}{\color{navy}\large\bfseries}{}{0pt}{}
\titlespacing{\section}{0pt}{11pt}{3pt}
\titleformat{\subsection}{\color{navy}\bfseries}{}{0pt}{}
\titlespacing{\subsection}{0pt}{7pt}{1.5pt}

\newtcolorbox{these}{colback=boxgray,colframe=boxedge,boxrule=0.5pt,arc=2.5pt,
  left=8pt,right=8pt,top=4.5pt,bottom=4.5pt}
\newtcolorbox{alerte}{colback=boxred,colframe=red!55,boxrule=0.8pt,arc=2.5pt,
  left=8pt,right=8pt,top=4.5pt,bottom=4.5pt}
\newtcolorbox{acquis}{colback=boxgreen,colframe=green!45,boxrule=0.8pt,arc=2.5pt,
  left=8pt,right=8pt,top=4.5pt,bottom=4.5pt}
\newtcolorbox{piste}{colback=boxplum,colframe=plum!45,boxrule=0.8pt,arc=2.5pt,
  left=8pt,right=8pt,top=4.5pt,bottom=4.5pt}
\newtcolorbox{extrait}{colback=boxcode,colframe=boxedge,boxrule=0.5pt,arc=1.5pt,
  left=7pt,right=7pt,top=3.5pt,bottom=3.5pt}

% \textcolor plutôt que {\color …} : en tête de cellule p{}, \color insère un whatsit
% qui ouvre une ligne vide et décale verticalement toute la colonne.
\newcommand{\best}[1]{\textcolor{green}{\textbf{#1}}}
\newcommand{\pire}[1]{\textcolor{red}{\textbf{#1}}}
\newcommand{\sz}[1]{\textcolor{grey}{\small\textsf{#1}}}
\newcommand{\code}[1]{\texttt{\small #1}}
\newcommand{\x}{$\times$}
\newcommand{\loc}[1]{{\footnotesize\color{grey}\texttt{#1}}}

\begin{document}

\begin{center}
{\Huge\bfseries\color{navy} Scella}\\[4pt]
{\large\color{grey}\itshape Le groupe le plus rachetable, maintenant}\\[4pt]
{\small\color{grey} 11 août 2026 \; · \; axe acquisition \; · \; les problèmes ne sont pas postulés : ils sont écrits dans le code}
\end{center}
\vspace{2pt}

\begin{these}
\textbf{Le résultat central, et il contredit mes deux rapports précédents.} Le groupe le plus
rachetable est \textbf{la pile preuve} -- huit crates, 12 862 lignes, 240 tests -- que le rapport
« produit unique » avait \textbf{explicitement jetée}, au motif qu'elle n'avait pas de producteur.
\textbf{C'est faux} : \code{provenance-capture} produit un DAG scellé sur le runtime, via quatre taps
\textbf{génériques} qui observent n'importe quel système d'IA, pas le moteur de ce dépôt. Elle est donc
la seule pièce dont la valeur \textbf{n'exige d'adopter aucune thèse} -- et la seule dont le problème
soit \textbf{écrit dans la loi} : \code{reg-evidence} nomme l'EU AI Act art.\,12, la CSDDD/LkSG et le
GDPR. À l'inverse, l'essaim réfutable -- vainqueur des deux analyses précédentes -- \textbf{tombe
dernier}, parce qu'une chose peut être la plus révolutionnaire \emph{et} la moins rachetable.
\end{these}

\section{1 · Ce que « rachetable maintenant » exige, et la contrainte que je maintiens}

« Rachetable » suppose un acheteur ; « répondre à un problème » suppose quelqu'un qui souffre. L'axe
est donc, pour la première fois de la série, tourné vers l'extérieur.

\begin{alerte}
\textbf{Je maintiens néanmoins la contrainte constante : je ne nomme aucun acquéreur et n'invente aucun
marché.} Ce n'est pas une précaution rhétorique -- c'est devenu inutile, pour une raison que la mesure
a révélée. \textbf{Les problèmes sont écrits dans le code}, article de loi à l'appui. Je les
\emph{lis}, je ne les \emph{postule} pas. Un marché inventé aurait été une opinion ; un docstring qui
cite « EU AI Act art.\,12 » est un fait vérifiable en une commande.
\end{alerte}

Cinq critères, choisis pour être évaluables sans destinataire nommé :

\begin{enumerate}[leftmargin=1.5em,itemsep=1.5pt,topsep=2pt]
\item \textbf{Problème daté et forcé} -- existe-t-il aujourd'hui hors du dépôt, avec une contrainte qui
      le rend non optionnel ? Non pas « comptera », mais \textbf{coûte déjà}.
\item \textbf{Transférabilité} -- tourne-t-il dans la pile de quelqu'un d'autre, \textbf{sans que
      l'acheteur adopte les thèses du dépôt} ?
\item \textbf{Preuve en une après-midi} -- une équipe de diligence peut-elle vérifier la revendication
      centrale \textbf{sans moi} ?
\item \textbf{Indépendance à une croyance} -- la valeur exige-t-elle d'accepter un pari ?
\item \textbf{Propreté juridique} -- licence, corps étrangers, nommage.
\end{enumerate}

Deux remarques de méthode. Le critère 3 est précisément celui que l'analyse « révolutionnaire » avait
\emph{déclassé} ; il redevient central ici. Et le critère 4 est celui qui décide -- il n'apparaissait
dans aucun des trois axes précédents. \textbf{Attendez-vous donc à un classement différent : c'est le
signe que l'axe est réellement nouveau, et non un habillage du verdict de la veille.}

\section{2 · Le fait qui renverse tout : la pile preuve a son \emph{propre} producteur}

\begin{alerte}
\textbf{Une élimination que j'ai faite hier était fausse, et il faut le dire avant tout le reste.} Dans
le rapport « produit unique », j'ai écarté la pile preuve d'un raisonnement en apparence solide :
\emph{sans producteur, elle n'a rien à sceller -- c'est une bibliothèque}. Je ne l'avais pas vérifié.
Cette fois, si.
\end{alerte}

\code{provenance-capture} l'annonce en en-tête : \emph{« Produit un \textbf{vrai DAG scellé} sur le
runtime »}. Il porte \textbf{quatre adaptateurs d'observation} -- appel modèle, message inter-agents,
consultation outil/retrieval, revue humaine -- qui convertissent des événements observés en nœuds
scellables. Il hache les entrées sensibles (minimisation GDPR), porte le cas dans l'identifiant du
nœud (donc scellé \emph{et} requêtable), embarque son \textbf{propre garde de souveraineté} -- egress
refusé, loopback uniquement, records persistés localement -- et \textbf{ne dépend que de
\code{provenance-dag}}, ce qui garde la chaîne de confiance mince.

\begin{acquis}
\textbf{Ces quatre taps sont génériques.} Ils n'observent pas le moteur de ce dépôt : ils observent
\emph{n'importe quel} système d'IA -- un appel de modèle, un message entre agents, une consultation
d'outil, une revue humaine sont les quatre événements de toute architecture agentique. C'est la
propriété la plus importante de toute cette analyse, parce que c'est elle qui rend l'actif transférable
sans adoption.
\end{acquis}

\section{3 · Les problèmes -- nommés par le code, article par article}

\code{reg-evidence} ne demande pas qu'on lui invente un marché ; il écrit le sien, et je le cite plutôt
que de le paraphraser :

\begin{extrait}
\sz{« transformer le DAG scellé re-prouvable en l'\textbf{artefact précis} qu'une équipe conformité doit
produire \textbf{de toute façon} --- registre de tenue de logs (\textbf{EU AI Act art.\,12}), dossier de
diligence fournisseur (\textbf{CSDDD/LkSG}), registre de traitement (\textbf{GDPR}) --- et le produire
\emph{mieux} que la main : automatiquement, traçablement, et \textbf{re-prouvable par le destinataire
lui-même}. »}
\end{extrait}

Trois obligations réglementaires \textbf{distinctes}, dont une datée par la loi. C'est la réponse la
plus directe possible à la question posée : non pas « un problème », mais \textbf{trois, nommés, non
optionnels}.

Et le gain est énoncé dans les termes mêmes qu'une équipe conformité emploie -- \textbf{heures
$\rightarrow$ minutes} -- sur un livrable qu'elle produit \emph{déjà}, à la main. C'est la forme de
valeur la plus facile à faire admettre : elle ne crée pas une dépense, elle en remplace une.

\section{4 · Les quatre candidats, mesurés}

\renewcommand{\arraystretch}{1.24}
\begin{center}\small
\begin{tabular}{@{}c >{\raggedright\arraybackslash}p{5.35cm} c c c >{\raggedright\arraybackslash}p{5.0cm}@{}}
\toprule
& \textbf{Le groupe} & \sz{Crates} & \sz{Lignes} & \sz{Tests} & \textbf{Les problèmes adressés} \\
\midrule
\textbf{V1} & \textbf{La preuve réglementaire} \sz{capture, DAG, verify, API, console, export forensique, gabarits, vérifieur 0-dép} & \best{8} & \best{12 862} & 240 & \textbf{AI Act art.\,12} · \textbf{CSDDD/LkSG} · \textbf{GDPR} · non-fabrication \\
\textbf{V2} & \textbf{La souveraineté d'egress} \sz{boot guard fail-closed, allowlist + cap, aveu de clair, \code{pii.rs}, reçu par saut, route bundle + ancre} & 11 & --- & --- & Fuite de périmètre en secteur régulé · attribution de coût LLM \\
\textbf{V3} & \textbf{L'essaim réfutable} \sz{le groupe couronné par mes deux analyses précédentes} & 10 & --- & \best{515} & Attribution multi-agents \sz{(émergent)} \\
\textbf{V4} & \textbf{Le vérifieur seul} \sz{\code{verify-mini} / \code{verify-meta-mini}} & \best{1} & --- & 29 & « Je ne peux pas auditer l'outil d'audit » \\
\bottomrule
\end{tabular}
\end{center}

\section{5 · Le classement}

Rangs de 1 à 4 par critère.

\renewcommand{\arraystretch}{1.2}
\begin{center}\small
\begin{tabular}{@{}r >{\raggedright\arraybackslash}p{6.0cm} ccccc c@{}}
\toprule
& \textbf{Le groupe} & \sz{Pb forcé} & \sz{Transfér.} & \sz{Preuve} & \sz{Sans croyance} & \sz{Juridique} & \textbf{Moy} \\
\midrule
\textbf{1} & \textbf{V1 · La preuve réglementaire} & \best{1} & \best{1} & 2 & \best{1} & \best{1} & \best{1,20} \\
2= & V2 · La souveraineté d'egress & 2 & 3 & \pire{4} & 2 & 3 & \textbf{2,80} \\
2= & V4 · Le vérifieur seul & \pire{4} & 2 & 3 & 3 & 2 & \textbf{2,80} \\
\textbf{4} & \textbf{V3 · L'essaim réfutable} & 3 & \pire{4} & \best{1} & \pire{4} & \pire{4} & \pire{3,20} \\
\bottomrule
\end{tabular}
\end{center}

\textbf{V1 est premier sur quatre critères sur cinq}, et l'écart -- 1,20 contre 2,80 -- est le plus
large de toute la série. Les raisons se cumulent au lieu de se compenser, ce qui est rare et mérite
d'être détaillé.

\textbf{Il ne demande aucune adoption.} Un acheteur branche quatre taps sur \emph{son} système ; il
n'adopte ni \code{efe-core}, ni l'essaim, ni aucune thèse sur les modèles latents. C'est la seule pièce
du dépôt dont la valeur ne dépend d'aucun pari -- et un pari est exactement ce qu'une revue
d'acquisition cherche à éliminer.

\textbf{Son problème porte une échéance légale.} Aucun des trois autres ne peut en dire autant. Un
problème réglementaire ne se négocie pas : il a une date, un destinataire et une sanction.

\textbf{Il est juridiquement propre.} Ses huit crates sont uniformément \code{UNLICENSED}. Les
\textbf{quatre seuls fichiers} du dépôt portant Apache-2.0 ou MIT sont ailleurs -- trois dans les tests
\code{m1d\_*} et un dans \code{efe-nano-sdk} -- c'est-à-dire \textbf{tous dans le périmètre de V3}.

\section{6 · Le renversement : mes deux analyses précédentes avaient écarté le meilleur actif}

Il faut le dire sans détour, et le tableau le dit mieux qu'un paragraphe.

\renewcommand{\arraystretch}{1.18}
\begin{center}\small
\begin{tabular}{@{}>{\raggedright\arraybackslash}p{4.6cm} >{\raggedright\arraybackslash}p{11.2cm}@{}}
\toprule
\textbf{Analyse} & \textbf{Verdict porté sur V1 (la pile preuve)} \\
\midrule
Ensembles \sz{11 août} & 2\textsuperscript{e} --- « la boucle pile B », classée comme une greffe \\
\textbf{Produit unique} & \pire{explicitement jetée} --- listée parmi les 57 crates à sortir du produit \\
Révolutionnaire & présente seulement par \code{forensic-export}, 1 levée sur 6 \\
\textbf{Rachetable} & \best{1\textsuperscript{re}, et de loin} \\
\bottomrule
\end{tabular}
\end{center}

D'où vient l'écart ? Mes critères précédents privilégiaient \textbf{la démonstrabilité et l'unicité de
la promesse}, et V1 se démontre plus lourdement qu'en une commande. Aucun de ces critères ne mesurait
ce qui décide ici : \emph{quelqu'un a-t-il déjà l'obligation de résoudre ce problème ?} Ce n'est pas
que les analyses précédentes étaient mal faites -- c'est que la question n'était pas la même, et qu'un
axe non mesuré ne se devine pas.

\textbf{Et V3, vainqueur des deux analyses précédentes, tombe dernier.} Non parce qu'il est faible : il
garde la meilleure démonstration du dépôt, seul critère où il reste premier. Mais sa valeur
\textbf{exige d'accepter une thèse} -- l'essaim latent, le méta-consensus --, il est le moins
transférable des quatre, et il porte à lui seul les quatre fichiers à licence tierce.

\begin{piste}
Une chose peut être la \textbf{plus révolutionnaire} et la \textbf{moins rachetable}. Ce n'est pas une
contradiction entre mes rapports : ce sont deux questions différentes, et le mérite de les avoir posées
séparément est précisément qu'on peut lire cet écart au lieu de le subir.
\end{piste}

\section{7 · Ce qu'une revue diligente verrait en premier}

Quatre points vérifiés, bloquants à des degrés divers, tous bon marché à corriger. Les énumérer
maintenant vaut mieux que de les laisser trouver.

\textbf{\code{main} est figée au 12 juillet et accuse 36 commits de retard.} Une équipe de diligence
clone \code{main} et ne voit \textbf{rien} de ce qui précède -- ni l'ancre, ni la liaison v4, ni le
producteur de bundle. C'est le blocage le plus grave, et le moins coûteux à lever.

\textbf{Le nom d'une autre entreprise est dans le code.} « Kovant » apparaît dans \textbf{23 fichiers
\code{.rs}}, dont \textbf{11 dans la pile preuve} -- jusqu'en en-tête de \code{provenance-capture}
(« Kovant P1 / S1 »). Techniquement anodin ; en lecture, un acquéreur ouvre le premier fichier de
l'actif convoité et y trouve le nom d'un tiers.

\textbf{Aucun profil de déploiement.} Établi par le rapport d'état des lieux : sept interrupteurs
d'environnement, et aucun fichier ni document ne les allume ensemble.

\textbf{Le garde de souveraineté est dupliqué de part et d'autre de la frontière} --
\code{provenance-capture/src/guard.rs} (dans V1) et \code{efe-swarm-sdk/src/backend/mod.rs} (dans V2).
Acquérir les deux modules impose donc de les fusionner, et les deux copies ont \textbf{déjà divergé une
fois} en produisant un contournement SSRF.

\section{8 · Le jugement}

\begin{acquis}
\textbf{Le groupe le plus rachetable maintenant est V1 -- la preuve réglementaire : huit crates,
12 862 lignes, 240 tests.}
\end{acquis}

Il est le seul du dépôt à réunir les trois conditions qui font une acquisition plutôt qu'une
curiosité : un \textbf{problème que quelqu'un est légalement obligé de résoudre}, une \textbf{valeur
qui ne dépend d'aucune croyance}, et une \textbf{vérification que l'acheteur mène lui-même} -- le
destinataire re-prouve l'artefact avec un binaire dont l'arbre de dépendances est vide, et l'export lui
dit de surcroît, dans le bundle, \emph{ce qu'il ne prouve pas}.

\textbf{Je ne recommande pas de le fusionner avec V2.} Les deux répondent à des problèmes voisins -- la
preuve et le périmètre -- mais V2 exige que l'acheteur fasse transiter son trafic par la passerelle,
c'est-à-dire un engagement d'architecture, quand V1 ne demande que quatre points d'observation.
Garder V2 comme \textbf{second module optionnel} préserve la propriété qui fait toute la valeur de V1 :
on peut l'adopter sans rien changer. Et si les deux partent ensemble, fusionner les gardes d'abord.

\textbf{Trois actions, par ordre de rendement :} fusionner la branche dans \code{main} · retirer
« Kovant » des 23 fichiers · écrire le profil.

\begin{these}
\small\textbf{Rappel de cadrage, inchangé dans toute la série.} Le reçu atteste l'\textbf{intégrité}
d'une suite de pas, jamais la \textbf{correction} d'une décision. Il est \emph{tamper-evident pour un
détenteur de clé}, pas \emph{trustless pour un tiers}. Les participants du mode réel exécutent une
fusion latente Mode-1 numérique : ce ne sont pas des « vrais SLM ». Et ce document évalue des
\emph{propriétés d'acquérabilité} -- il ne prédit aucun acquéreur et n'en nomme aucun.
\end{these}

\section{Annexe · Vérifier chaque affirmation soi-même}

Les \emph{notes} des cinq critères sont un jugement. Les \textbf{mesures} et les \textbf{citations},
elles, se rejouent -- et ce sont elles qui portent le verdict.

\renewcommand{\arraystretch}{1.22}
\begin{center}\small
\begin{tabular}{@{}>{\raggedright\arraybackslash}p{5.3cm} >{\raggedright\arraybackslash}p{10.8cm}@{}}
\toprule
\textbf{Affirmation} & \textbf{Ce qui la rejoue} \\
\midrule
V1 est de 8 crates &
  \code{ls crates | sort > /tmp/ws.txt} \newline
  \code{for p in reg-evidence forensic-export audit-api audit-console \textbackslash} \newline
  \code{\ \ provenance-capture provenance-verify verify-mini ; do \textbackslash} \newline
  \code{\ \ cargo tree -p \$p -e normal -{}-prefix none ; done \textbackslash} \newline
  \code{\ \ | awk '\{print \$1\}' | sort -u | comm -12 - /tmp/ws.txt} \\
La pile preuve \textbf{produit} &
  \code{sed -n '1,22p' crates/provenance-capture/src/lib.rs} \\
Les problèmes, par article &
  \code{sed -n '1,16p' crates/reg-evidence/src/lib.rs} \\
12 862 lignes, 240 tests &
  \code{find crates/\{provenance-dag,provenance-capture,…\} -name '*.rs' | xargs wc -l} \\
Licence : 4 fichiers non propriétaires, tous hors V1 &
  \code{grep -rl 'SPDX-License-Identifier: \textbackslash(Apache-2.0\textbackslash|MIT\textbackslash)' crates -{}-include='*.rs'} \\
« Kovant » dans 23 fichiers &
  \code{grep -rl 'Kovant' crates -{}-include='*.rs' | wc -l} \\
\code{main} accuse 36 commits &
  \code{git log -1 -{}-format=\%ci main \; \&\& \; git rev-list -{}-count main..HEAD} \\
Le garde est dupliqué &
  \code{grep -rln 'fn is\_local\_endpoint' crates -{}-include='*.rs'} \\
\bottomrule
\end{tabular}
\end{center}

Un avertissement pour finir, et il porte sur la nature de ce document. Les empreintes, les citations et
les zéros sont des \textbf{faits}, datés du 11 août. L'appréciation « ce problème est forcé » repose
sur un texte réglementaire que le code cite lui-même -- c'est plus solide qu'une opinion de marché,
mais ce n'est pas une mesure : la portée réelle d'une obligation dépend du secteur et de la taille de
celui qui y est soumis, et ce document n'en sait rien.

\end{document}
